% ---------------------------------------------------------------------------
% What the first-order DDP scheme exchanges, and how slowly. Drawn for T = 4
% stages over 4 sweeps -- the smallest grid that shows terminal information
% completing its walk from the last stage to the first.
%
% Reading the layout: ROWS are sweeps (one pass of the outer loop), COLUMNS are
% stages. One cell is one small QP. The two couplings are deliberately drawn at
% different angles because that difference IS the mechanism:
%   horizontal, solid   B[t-1] from THIS sweep -- a hard constraint, so it moves
%                       left to right within a row and is always current
%   diagonal, dashed    mu[t+1] and B[t+1] from the PREVIOUS sweep -- a soft
%                       linear term, so it can only step one column left per row
% The diagonal is why the terminal stage's information needs T sweeps to reach
% stage 1; the highlighted staircase traces exactly that walk.
%
% Colour semantics follow the tADMM exchange figure of the companion study, so
% the same narrative carries across both:
%   blue  = local / primal, owned by the sweep that solves it   (B[t])
%   red   = stale, owned by no sweep in progress                (B[t+1] frozen)
%   green = dual, the feedback that closes the loop             (mu[t+1])
%
% GREYSCALE PROOFING. Hue is never the only cue:
%   solved cell     solid dark fill, sharp corners, heavy border
%   stale parameter near-white fill, DASHED border
%   dual            DASHED line only, never a fill, so it cannot collide with
%                   red under deuteranopia
%   the staircase   drawn heavier than every other arrow
% ---------------------------------------------------------------------------
\definecolor{ddpLoc}{HTML}{1F5186}     % deep blue -- primal, solved this sweep
\definecolor{ddpLocE}{HTML}{143A61}
\definecolor{ddpStale}{HTML}{C0392B}   % red       -- frozen from last sweep
\definecolor{ddpStaleE}{HTML}{8E2A20}
\definecolor{ddpDual}{HTML}{0F6B45}    % green     -- the dual mu

\begin{figure*}[!t]
\centering
\begin{tikzpicture}[
  x=1mm, y=1mm, font=\small,
  cell/.style ={draw=ddpLocE, line width=0.5pt, minimum width=16mm,
                minimum height=7.6mm, inner sep=0pt, fill=ddpLoc, text=white},
  hard/.style ={-{Stealth[length=1.6mm,width=1.3mm]}, draw=ddpLocE,
                line width=0.7pt},
  soft/.style ={-{Stealth[length=1.6mm,width=1.3mm]}, draw=ddpDual,
                line width=0.5pt, densely dashed},
  stair/.style={-{Stealth[length=2.0mm,width=1.6mm]}, draw=ddpStaleE,
                line width=1.1pt, densely dashed},
  hdr/.style  ={font=\small\bfseries},
]
\def\cp{24}   % column pitch (stages)
\def\rp{14}   % row pitch    (sweeps)

% ---- headers ----
\foreach \t in {1,2,3,4}
  \node[hdr] at ({\cp*(\t-1)},{11}) {$t=\t$};
\node[hdr, anchor=south, rotate=90] at (-25,{-\rp*1.5}) {sweep $k$};

% ---- the grid of stage solves ----
\foreach \k in {1,2,3,4} {
  \node[hdr, anchor=east] at (-14,{-\rp*(\k-1)}) {$k=\k$};
  \foreach \t in {1,2,3,4}
    \node[cell] (c\k\t) at ({\cp*(\t-1)},{-\rp*(\k-1)}) {$B^{\k}[\t]$};
}

% ---- hard coupling: within a sweep, left to right ----
\foreach \k in {1,2,3,4}
  \foreach \t/\tn in {1/2, 2/3, 3/4}
    \draw[hard] (c\k\t.east) -- (c\k\tn.west);

% ---- soft coupling: from the previous sweep, one column left ----
\foreach \k/\kp in {2/1, 3/2, 4/3}
  \foreach \t/\tn in {1/2, 2/3, 3/4}
    \draw[soft] (c\kp\tn.south) -- (c\k\t.north);

% ---- the staircase: terminal information walking to stage 1 ----
\draw[stair] (c14.south) -- (c23.north);
\draw[stair] (c23.south) -- (c32.north);
\draw[stair] (c32.south) -- (c41.north);

% ---- legend: one horizontal row below the grid, so it cannot overrun the
%      two-column text width however the entries are worded ----
\def\ly{-62}
\draw[hard] (-13,{\ly}) -- (-2,{\ly});
\node[anchor=north west, align=left, text width=40mm] at (0,{\ly+3})
  {\textbf{hard}: $B[t\!-\!1]$ from \emph{this} sweep, as a constraint};
\draw[soft] (46,{\ly}) -- (57,{\ly});
\node[anchor=north west, align=left, text width=42mm] at (59,{\ly+3})
  {\textbf{soft}: $\mu[t\!+\!1]$, $B[t\!+\!1]$ \emph{frozen} from the last sweep,
   entering the objective linearly};
\draw[stair] (105,{\ly}) -- (116,{\ly});
\node[anchor=north west, align=left, text width=38mm] at (118,{\ly+3})
  {\textbf{the walk}: one terminal fact, reaching stage $1$ in sweep $T$};

\end{tikzpicture}
\caption{What the first-order scheme passes between its subproblems, drawn
for four stages over four sweeps.  Each cell is one small quadratic program.
Read one row at a time.  Within a sweep the state moves left to right along the
solid arrows and is always current --- stage $t$ is bound to the energy stage
$t\!-\!1$ has just produced, as a hard constraint.  The information travelling
the other way cannot do that.  Stage $t$ learns what the future is worth only
through $\mu[t\!+\!1]$ and $B[t\!+\!1]$ as they stood when the \emph{previous}
sweep finished, so that knowledge slides one column left per row and no faster.
The heavy staircase follows a single fact on its way through: what the terminal
stage discovers in sweep~$1$ does not reach stage~$1$ until sweep~$4$.  A
horizon of $T$ periods therefore needs about $T$ sweeps before the first stage
has heard from the last, and each sweep costs $T$ solver calls.}
\label{fig:userddp}
\end{figure*}
